\documentclass{article}
\usepackage{amsmath}
\begin{document}
\begin{flushright}
CSE 21 Fall 2026\\
Homework 1\\
Submitted: {\today}
\end{flushright}

\begin{center}
Names: David Washington, Judah Boyce, Osinachi Agaranna, Nick Young\\
TSN: 200484943, 200453537, 200455712, 200455733
\end{center}

\begin{enumerate}
    \item
    \begin{enumerate}
      \item Assume that for all $n \ge 1$ 
        \[
            2 + 5 + 8 + \cdots + (3n - 1) = \frac{n(3n + 1)}{2}
        \]
        First we prove that at $n=1$ the left and right side are equivalent: P(1) = P(1).  When $n = 1$:
        \[
          (3(1) - 1) 
          = 2
          = \frac{1(3(1) + 1)}{2}
          =  \frac{4}{2}
          = 2
        \]
        Assuming $P(k)$ is true. For $n \ge 1$ , $P(k) \rightarrow P(k+1)$ is again equivalent for the left and right side.
        RHS $P(k+1)$
        \[
          \frac{(k+1)(3(k+1) + 1)}{2}
          = \frac{(k+1)(3k+4)}{2}
          = \frac{3k^2+7k+4}{2}
        \]
        Since the LHS and RHS are equivalent we add the next element of the LHS to both sides and it should be equivalent to the RHS P(k+1)
        \[
          2 + 5 + 8 + \cdots + (3k - 1) + (3(k+1)-1)
          = \frac{k(3k+1)}{2} + (3(k+1)-1)
        \]
        \[
          = \frac{(k)(3k+1)}{2} + \frac{6k+4}{2} 
        \]
        \[
          = \frac{3k^2+k}{2} + \frac{6k+4}{2} 
        \]
        \[
          = \frac{3k^2+7k+4}{2}
        \]

      The last equation shows that $P(k+1)$ is true under the assumption that $P(k)$ is also true. Therefore,
        \[
            2 + 5 + 8 + \cdots + (3n - 1) = \frac{n(3n + 1)}{2}
        \]
        Assume that for all $n \ge 1$ 

      \item Assume for all $n \geq 1$:
      \[
        \prod_{k=1}^{n}\left(1 + \frac{1}{k}\right) = n + 1
      \]
      First we prove P(k) is true such that at $n=1$ the left and right side are equivalent:
      \[
        \prod_{k=1}^{1}\left(1 + \frac{1}{k}\right) = 1 + 1
      \]
      \[
        (1+1) = (1+1)
      \]
      \[
       2 = 2
      \]
      Since P(k) is true we will now prove that $P(k) \rightarrow P(k+1)$ is again equivalent for the left and right side.
      LHS:
      \[
        = (k+1) + 1
        =k+2
      \]
      RHS and LHS:
      \[
        \prod_{k=1}^{n}\left(1 + \frac{1}{k+1}\right) = (k+1)(1+\frac{1}{k+1})
      \]
      \[
       = k+1 + \frac{k+1}{k+1}
       = k + 1 + 1
       = k + 2
      \]

      This proves that $P(k+1)$ is true assuming that $P(k)$ is true.
      Therefore, by mathematical induction,
      \[
        \prod_{k=1}^{n}\left(1+\frac{1}{k}\right)=n+1
      \]
      for all $n\geq1$.

    \end{enumerate}
  
    \item A license plate is a string over the alphabet of 36 characters consisting of 26 uppercase Latin letters and 10 numerals (0 through 9). (2 points for a correct expression and 2 points for justification)
    \begin{enumerate}
      \item How many 7-character license plates have at least 2 numerals?

      The total number of combinations for 7 positions and 36 characters is $36^{7}$. To find the case where there is at least 2 we can find the compliment (number of combos with no numbers or one number) and subtract that from the total.
      \begin{itemize}
        \item Total number with no numbers: $26^{7}$
        \item Total number with one number: $\binom{7}{1} \cdot 10 \cdot 26^{6}$
      \end{itemize}
      Number of license plates that have at least 2 numerals:
      \[
         = 36^{7} - 26^{7} - \binom{7}{1} \cdot 10 \cdot 26^{6}
      \]

      \item How many 7-character license plates contain the word MILES as a consecutive substring?

      To find the possible combos, MILES is a 5 character string that must be in order in the license plate thus there are only 3 ways to place the specific string on the license plate. There are always 2 empty characters that allows for 36 characters each thus:
      \begin{itemize}
        \item Number of MILES placements: $3$
        \item Number of unique options for 2 spots with 36 characters: $36^{2}$
      \end{itemize}
      Number of plates that contain the word MILES as a consecutive substring:
      \[
        = 3 \cdot 36^{2}
      \]

      \item How many 7-character license plates consist of 3 different characters with 3 copies of
      one, 3 copies of another and a single copy of the last?
      Ex: AAA3331, P33T3TT, ...

      Choose 3 characters. Two of them appear 3 times, and the other appears once.
      \begin{itemize}
        \item Choose 3 characters from 36: $\binom{36}{3}$
        \item Choose which of those 3 appears once: $3$
        \item Arrange 7 spots with those repetitions: $\frac{7!}{3! \cdot 3! \cdot 1!}$
      \end{itemize}
      Number of ways to arrange 3 different characters with 3 copies of
      one, 3 copies of another and a single copy of the last:
      \[
        = \binom{36}{3} \cdot 3 \cdot \frac{7!}{3! \cdot 3! \cdot 1!}
      \]

      \item How many 7-character license plates have exactly 3 different numerals and 4 copies of
      the same letter?

      Similar to the previous one you are choosing 1 character but in this case an alphabetical character from 26 which is repeated 4 times also in this problem you have the last 3 characters as unique numerals
      \begin{itemize}
        \item Choose the letter from 26: $\binom{26}{1}$
        \item Choose and order 3 different numerals: $P(10,3)$
        \item Arrange 7 spots with that letter repeated 4 times: $\frac{7!}{4! \cdot 1!}$
      \end{itemize}
      The number of ways to have exactly 3 different numerals and 4 copies of
      the same letter:
      \[
         = \binom{26}{1} \cdot P(10,3) \cdot \frac{7!}{4! \cdot 1!}
      \]

      \item How many 7-character license plates are there that are all letters and they are in
      alphabetical order?

      For this one you choose 7 alphabetical characters from 26 and that set has one alphabetical order.      
      \begin{itemize}
        \item Choose the 7 from 26: $\binom{26}{7}$
      \end{itemize}
      The number of ways that the license plate is all letters and they are in
      alphabetical order:
      \[
         = \binom{26}{7}
      \]
    \end{enumerate}

    An art museum with 12 different rooms just got a shipment of 9 new paintings
    \begin{enumerate}
      \item How many ways can the museum curator arrange all 9 paintings among the 12 rooms
      in the museum?
      You can arrange 9 paintings and have each painting have 12 options to choose form when being placed.
      Thus the number of ways to place 9 paintings amon 12 rooms in the museum is:
      \[
         = 12^{9}
      \]

      \item How many ways can the museum curator arrange all 9 paintings among the 12 rooms
in the museum so that each painting is in a different room?
      This question asks the number of permutations of 9 paintings in 12 rooms.
     
      Thus the number of ways to place 9 paintings amon 12 rooms in the museum is:
      \[
         = P(12,9)
      \] 

      \item How many ways can the museum curator arrange all 9 paintings among the 12 rooms
      in the museum so that at least one painting is in the biggest room and exactly one painting is in
      the smallest room?
      For this question we consider the one painting chosen for the biggest room and that one painting must be in the smallest room.
      \begin{itemize}
        \item Number of ways to chosen 1 painting for the biggest room = $ \binom{9}{1} $
        \item Number of ways to place 8 paintings and at least one in the smallest room (compliment, total arrangements subtracted by possible arrangements not in the smallest room) = $11^8 - 10^8$ 
      \end{itemize}
      So the total number of arrangements for 9 paintings one in the biggest room and at least one in the smallest is:
      \[
        =  \binom{9}{1} \cdot 11^8 - 10^8
      \]

      \item Paintings A,B, and C are all painted by the same artist. How many ways can the
      museum curator arrange all 9 paintings among the 12 rooms such that paintings A,B, and C are
      all in the same room?
      We can consider paintings A, B, and C to be one set or basically one painting and so now we have 7 paintings being placed in 12 rooms
      So the total number of arrangements for all 9 paintings when A,B, and C must be in the same room is:
      \[
        =  12^{7}
      \]
      \item How many ways can the museum curator arrange all 9 paintings in 3 of the 12 rooms
      such that each room has exactly 3 paintings?
      Given that each room has exactly 3 paintings it means that only 4 rooms will be used at a time, so we can think of it in the perspective of those 4 rooms.
      \begin{itemize}
        \item For the first room you have 12 paintings and you need to choose 3 = $\binom{12}{3}$
        \item For the second room you have 9 paintings now and need 3 = $\binom{9}{3}$
        \item Same thing for the third =  $\binom{6}{3}$
        \item Same thing for the fourth  =  $\binom{3}{3}$
      \end{itemize}
      Thus the total number of ways is equal to product of each room:
      \[ 
        = \binom{12}{3} \cdot \binom{9}{3} \cdot \binom{6}{3} \cdot \binom{3}{3}
      \]
    \end{enumerate}
    \begin{enumerate}
    \item How many different ways are there to arrange the letters in the following string:
    INCLUSIONEXCLUSION? 
    For this problem the substring CLUSION is repeated twice and the letter I and N once more giving the following:
    \begin{itemize}
      \item Total number of arrangements = $18!$
      \item I and N repeated 3 times and C, L, U, S, O twice = $3! \cdot 3! \cdot 5(2!)$ 
    \end{itemize} 
    Different ways to arrange letters from the string INCLUSIONEXCLUSION:
    \[
    = \frac{18!}{3! \cdot 3! \cdot 5(2!)}
    \]

    \item How many different ways are there to color the 9 edges of a triangular prism with 9 different
    colors (so that two colorings are considered the same if one can be oriented to look the same as
    the other one?)
    For this one we consider the number of edges and colors for each and then the ways to rotate the shape
    \begin{itemize}
      \item Number of ways to place 9 colors on 9 edges $=9!$
      \item Number of ways to rotate the shape left and right and twist = $ 3 \cdot 2 $
    \end{itemize} 
    So the total number of ways to color the edges given same orientation:
    \[
    = \frac{9!}{2 \cdot 3}
    \]

    \item There are 18 different basketball players. How many ways can they organize themselves into 3
    games each with 3 players versus 3 players?
    For this one you have 3 different games with 6 people and then within those 3 games you have two teams that you need to place the 3 people on each team.
    \begin{itemize}
      \item First game: $ \binom{18}{6} $
      \item Second game: $ \binom{12}{6} $
      \item Third game: $ \binom{6}{6} $
      \item Ways to arrange the 3 teams for 3 people on two teams : $ 3 \binom{6}{3} $
    \end{itemize}
    There are The number of ways organize 18 people into 3
    games each with 3 players versus 3 players:
    \[
    = \binom{18}{6} \cdot \binom{12}{6} \cdot \binom{6}{6} \cdot 3 \binom{6}{3} 
    \]
    \end{enumerate}


 \item
 \begin{enumerate}

 \item How many different ways are there to arrange the letters in the following string:
 \[ 
 \text{INCLUSIONEXCLUSION}
 \]

 The string contains 18 letters total, with a repeat of letters with the following multiplicities:
 \[
 I:3, \quad N:3, \quad C:2, \quad L:2, \quad U:2, \quad S:2, \quad O:2.
 \]

 The letters $E$ and $X$ each occur once. Since rearranging identical copies of the same letter not produce a new arrangement we must divide by the factorial of the repeats of each letter:

 \[ 
 \boxed{ 
 \frac{18!}{3!3!(2!)^5}
 }
 \]

 \textbf{Justification:}
 We divide by $3!$ twice for the three copies of $I$ and $N$, and by $2!$ for each of the five letters that appear twice. This removes all duplicate arrangments.

 \item How many different ways are there to color the 9 edges of a triangular prism with 9 different colors, where two colorings are considered the same if one can be oriented to look the same as the other?

 If we first look at the amount of ways we could assign the 9 distinct colors to be assigned to the 9 distict edges we would get 
 \[
 9!
 \]
 ways.

 A triangular prism has 6 rotational symmetries and since there are 9 different colors no non-identiy rotation can leave a coloring unchanged. As a result, each distinct coloring is counted 6 times among the arrangements:

 \[ 
 \boxed{
 \frac {9!}{6}
 }
 \]

 \textbf{Justification:}
 The 6 possible rotations of the triangular prism produce colorings that are considered equivalent and because every edge has a different color, each group contains 6 rotated versions.

 \item There are 18 different basketball players. How many ways can they organize themselves into 3 games, each with 3 players versus 3 players?

 
 \[
 \boxed{
 \frac{18!}{(3!)^6 2^3 3!}
 }
 \]

\textbf{Justification:}
The factor $(3!)^6$ accounts for the fact that the order of the 3 players within each team does not matter. The factor $2^3$ accounts for swapping the two teams in each game. $3!$ accounts for rearranging the order of the three games .

\end{enumerate}

\end {enumerate}
 


\end{document}
